\chapter{支持向量回归的误差容忍带}
\label{chap:svr}
\begin{quote}\small
\textit{本章摘要。}支持向量回归允许小误差不产生损失，超过容忍带后按超出量线性计费。本章从损失定义推导松弛变量、拉格朗日对偶与支持向量，解释容忍带宽度和正则参数分别控制什么。
\end{quote}

第\ref{chap:krr}章的核岭回归用平方损失处理每个残差，残差越大，损失对预测值的梯度绝对值越大。若温度传感器偶有异常读数，这类样本可能显著影响拟合。支持向量回归（Support Vector Regression，SVR）采用不同的代价规则。它允许预测值与观测值在一定范围内不同，只有超出容忍范围的部分才产生损失。本章选自Theodoridis第11.8节\cite{theodoridis2015ml}，与核方法网页的SVR公式对照\cite{zheng2021kernel}。

\section{误差容忍带与不敏感损失}
给定监督回归训练集$\{(\bm x_i,y_i)\}_{i=1}^n$，其中$\bm x_i\in\R^d$、$y_i\in\R$，模型输出实数$f(\bm x_i)$。SVR适用于允许一定预测偏差而希望模型对更大偏差按线性代价处理的回归任务。令残差$r_i=y_i-f(\bm x_i)$，定义$\epsilon$不敏感损失（$\epsilon$-insensitive loss）
\begin{equation}
 L_\epsilon(r)=\max(0,|r|-\epsilon),\qquad\epsilon\ge0.
 \label{eq:svr-insensitive-loss}
\end{equation}
$\epsilon$与标签有相同单位。若标签是摄氏温度，$\epsilon=0.5$表示误差不超过$0.5$摄氏度时训练损失为零。残差为$0.2,0.8,-1.2$时，损失分别为$0,0.3,0.7$。

当$r>\epsilon$时，损失对残差的斜率为1。当$r< -\epsilon$时，斜率为$-1$。在边界点损失不可微，可用次梯度描述。斜率的绝对值不随误差增大，因此其尾部响应小于平方损失，但损失仍随极端误差线性增长。Huber损失则在小误差处为二次形式，在大误差处转为线性形式。$\epsilon$带只规定训练损失为零的残差范围，不提供概率覆盖率保证，因此不是置信区间。

图\ref{fig:svr-tube}给出具体曲线$f(x)=x$及$\epsilon=0.5$的带。上下虚线分别是$f+\epsilon$与$f-\epsilon$，松弛量只计算带外部分。
\begin{figure}[htbp]
\centering
\begin{tikzpicture}[scale=0.9,>=Latex]
 \draw[->] (-0.3,0)--(4,0) node[right] {$x$};
 \draw[->] (0,-0.4)--(0,4.1) node[above] {$y$};
 \fill[blue!7] (0,0.5)--(3.2,3.7)--(3.2,2.7)--(0,-0.5)--cycle;
 \draw[thick] (0,0)--(3.3,3.3) node[right] {$f(x)=x$};
 \draw[dashed] (0,0.5)--(3.2,3.7) node[right] {$f+\epsilon$};
 \draw[dashed] (0,-0.5)--(3.2,2.7) node[right] {$f-\epsilon$};
 \fill (1,1.2) circle (0.06);
 \fill (2,2.8) circle (0.06);
 \draw[red,thick] (2,2.5)--(2,2.8) node[midway,right] {$\xi=0.3$};
 \node[below right] at (1,1.2) {带内点};
\end{tikzpicture}
\caption{教学示意的$\epsilon$容忍带，不是实验拟合结果。观测$(1,1.2)$在带内，$(2,2.8)$比带上缘高0.3，因此上侧松弛量为0.3。}
\label{fig:svr-tube}
\end{figure}

\section{将损失改写为约束优化}
用$f(\bm x)=\langle\bm w,\phi(\bm x)\rangle+b$预测实数，其中$\phi(\bm x)$是特征映射，$\bm w$属于相应特征空间，$b\in\R$是不受惩罚的截距。引入非负松弛量$\xi_i,\xi_i^*$，分别度量观测高于上缘和低于下缘的超出量：
\begin{align}
 \min_{\bm w,b,\bm\xi,\bm\xi^*}\quad
 &\tfrac12\|\bm w\|_{\mathcal H}^2+C\sum_i(\xi_i+\xi_i^*),\label{eq:svr-primal}\\
 y_i-f_i&\le\epsilon+\xi_i,\qquad
 f_i-y_i\le\epsilon+\xi_i^*,\quad \xi_i,\xi_i^*\ge0.\nonumber
\end{align}
这里$f_i=f(\bm x_i)$，$C>0$是带外误差的惩罚系数。第一条约束等价于$\xi_i\ge r_i-\epsilon$，第二条等价于$\xi_i^*\ge-r_i-\epsilon$。目标对两类松弛量都递增，因此固定函数时最优取值为
\[
 \xi_i=\max(0,r_i-\epsilon),\qquad
 \xi_i^*=\max(0,-r_i-\epsilon).
\]
两者之和恰为式\eqref{eq:svr-insensitive-loss}中的单样本损失。故约束优化等价于$\tfrac12\|\bm w\|_{\mathcal H}^2+C\sum_{i=1}^nL_\epsilon(r_i)$。此处损失求和而非平均，若改写为平均损失，正则项与$C$的相对尺度必须同步换算。

\section{正负展开系数与对偶问题}
对两条带约束分别引入$\alpha_i,\alpha_i^*\ge0$，对非负松弛量引入$\beta_i,\beta_i^*\ge0$。拉格朗日函数为
\begin{align}
 \mathcal L={}&\tfrac12\|\bm w\|^2+C\sum_i(\xi_i+\xi_i^*)
 +\sum_i\alpha_i(y_i-f_i-\epsilon-\xi_i)\nonumber\\
 &+\sum_i\alpha_i^*(f_i-y_i-\epsilon-\xi_i^*)
 -\sum_i(\beta_i\xi_i+\beta_i^*\xi_i^*).
\end{align}
对$\bm w$求导得到特征空间驻点条件。对$b$求导得到等式约束。对两类松弛量求导得到乘子上界：
\begin{equation}
 \bm w=\sum_{i=1}^n(\alpha_i-\alpha_i^*)\phi(\bm x_i),\qquad
 \sum_{i=1}^n(\alpha_i-\alpha_i^*)=0,
 \quad C-\alpha_i-\beta_i=0,\quad C-\alpha_i^*-\beta_i^*=0.
\end{equation}
因$\beta_i,\beta_i^*\ge0$，故$0\le\alpha_i,\alpha_i^*\le C$。记带符号系数$\bm a=\bm\alpha-\bm\alpha^*$，用$K_{ij}=\langle\phi(\bm x_i),\phi(\bm x_j)\rangle=k(\bm x_i,\bm x_j)$表示核矩阵。将驻点代回拉格朗日函数，得到
\begin{equation}
 \max_{\bm\alpha,\bm\alpha^*}
 -\tfrac12\bm a^{\mathsf T}K\bm a+\bm y^{\mathsf T}\bm a
 -\epsilon\bm1^{\mathsf T}(\bm\alpha+\bm\alpha^*),
 \quad\bm1^{\mathsf T}\bm a=0,
 \quad0\le\alpha_i,\alpha_i^*\le C.
 \label{eq:svr-dual}
\end{equation}
这里$K$是半正定训练核矩阵。原问题是凸优化问题。取$\bm w=0,b=0$并选择足够大的正松弛量，可使所有不等式约束严格成立，满足Slater条件，因此原始最优值与对偶最优值相等。预测为
\begin{equation}
 \hat y(\bm x)=\sum_i a_i k(\bm x_i,\bm x)+b.
 \label{eq:svr-prediction}
\end{equation}
与分类SVM的展开系数$\alpha_i y_i$不同，SVR的展开系数是$\alpha_i-\alpha_i^*$。标签已通过对偶目标和约束影响这些乘子，不应再乘入式\eqref{eq:svr-prediction}。

\section{支持向量与截距的计算}
KKT条件由原始可行性、对偶可行性、驻点条件和互补松弛组成。对偶变量与约束的互补关系为
\begin{align}
 \alpha_i(r_i-\epsilon-\xi_i)&=0,
 &(C-\alpha_i)\xi_i&=0,\\
 \alpha_i^*(-r_i-\epsilon-\xi_i^*)&=0,
 &(C-\alpha_i^*)\xi_i^*&=0.
 \label{eq:svr-kkt-complementarity}
\end{align}
严格位于带内时，两条带约束均有余量，故$\alpha_i=\alpha_i^*=0$，该点不进入预测展开。若$0<\alpha_i<C$，第二个互补条件给出$\xi_i=0$，第一个条件随即给出$r_i=\epsilon$，即点在上缘。若$0<\alpha_i^*<C$，同理有$r_i=-\epsilon$，点在下缘。若$r_i>\epsilon$，则上侧松弛量$\xi_i>0$且$\alpha_i=C$。若$r_i< -\epsilon$，则下侧松弛量$\xi_i^*>0$且$\alpha_i^*=C$。当$\epsilon>0$时，同一样本不可能同时处于带的两侧，因此$\alpha_i$与$\alpha_i^*$不会同时为正。带边界上的点也可能有零系数，实际预测只由非零系数样本构成。

对上缘自由支持向量$s$与下缘自由支持向量$t$，分别有
\begin{equation}
 b=y_s-\epsilon-\sum_i a_iK_{is},\qquad
 b=y_t+\epsilon-\sum_i a_iK_{it}.
\end{equation}
其中$K_{is}=k(\bm x_i,\bm x_s)$。对多个自由支持向量，可分别计算截距后取平均，以减小数值误差。若没有自由支持向量，可由KKT不等式确定满足条件的截距区间。具体而言，记$g_i=\sum_j a_jK_{ji}$。若$\alpha_i=0$，则$y_i-g_i\le b+\epsilon$。若$\alpha_i=C$，则$b+\epsilon\le y_i-g_i$。若$\alpha_i^*=0$，则$b-\epsilon\le y_i-g_i$。若$\alpha_i^*=C$，则$y_i-g_i\le b-\epsilon$。这些上下界共同约束$b$，数值解应在容差内满足其交集。退化情形下，恰位于带边界的样本可以有零系数，因此支持向量以非零对偶系数为准。

\subsection{两点线性算例}
取一维线性模型$f(x)=wx+b$，训练点为$(x_1,y_1)=(-1,-1)$和$(x_2,y_2)=(1,1)$，并令$\epsilon=0.2$。数据关于原点对称，最优解可取$b=0$。由此目标化为
\begin{equation}
 J(w)=\tfrac12w^2+2C\max(0,|1-w|-0.2).
\end{equation}
在$0\le w<0.8$区间，导数为$w-2C$。在$0.8\le w\le1.2$区间，损失为零，范数最小点是$w=0.8$。其他区间不能得到更低值，因此$w=\min(2C,0.8)$。

当$C>0.4$时，$w=0.8$，两点的残差分别为$-0.2$和$0.2$，都落在带边界。由驻点条件$w=a_2-a_1$和等式约束$a_1+a_2=0$，得$a_1=-0.4,a_2=0.4$，对应的非零乘子为$0.4<C$。当$C=0.1$时，$w=0.2$，两点残差绝对值均为$0.8$，各自松弛量为$0.6$，相应非零乘子达到上界$C$。在$C=0.4$时，最优斜率仍为$0.8$，乘子恰好达到上界。这个算例显示，较强的范数惩罚会以更大的带外残差为代价。

\section{参数选择与适用边界}
输入缩放会改变核距离，输出缩放会改变$\epsilon$的量纲和数值含义。缩放参数只能在训练折上估计，再用于相应验证数据。应联合验证$C$、$\epsilon$与核尺度。增大$\epsilon$会扩大零损失区域，但最优函数和系数也会随之改变，因此支持向量数不一定单调下降。单凭带宽不能保证指定稀疏率。

SVR减少了大误差的梯度增长，但不能防止错误标签、分布变化或高影响输入位置造成偏差。它输出点预测，不提供预测区间。若目标是条件分位数或带概率覆盖率的区间，应使用对应损失或另行校准。测试时要同时检验误差、异常读数影响和所需预测成本。

\section{阅读依据}
主要依据为Theodoridis第11.8节（书内页530--537），正文固定$\xi_i$处理上侧残差的符号约定，使原始约束、对偶、截距和算例一致。
